\providecommand{\myrootdir}{..}
\documentclass[\myrootdir/main.tex]{subfiles}

\begin{document}

\chapter{Helpful snippets}

This chapter just randomly uses some snippets and frequently used code patterns you probably might find useful.
The content is pure nonsense, but the underlying LaTeX code can easily be modified and adapted to the real text of your thesis.

\section{Structuring text}

First of all, there are a few elements that help you structure your text.

\subsection{Lists and enumerations}

Sometimes using lists and enumerations are the best way to present your data, although some authors see this as a bad writing style.

\begin{itemize}
 \item One item
 \item Another item
 \item the last item
\end{itemize}

\noindent Or enumerations are also possible

\begin{enumerate}
 \item Number one
 \item here comes number two
 \item All good things come in threes
\end{enumerate}

\subsection{Fancy lists}

But if you want to use lists, there are also some ways to get a better layout with them.

\begin{enumerate}[label=\textbf{[\Roman*]}]
 \item Even more enumerations
 \item But this time with custom style
 \item \ldots this list also needs a third item
 \item and a fourth to show roman enumerations better \ldots
\end{enumerate}

Or just use them for better alignment

\begin{description}[align=right,labelwidth=2cm]
 \item [Head] \texttt{~(Q.Q)~}
 \item [Body] \texttt{(/~n~\textbackslash{})}
 \item [Feet] \texttt{~~"~"~~}
 \item [Info] In the rare case you do not recognize it, this is an ASCII-owl.
\end{description}


\section{Using pictures}

Showing some pictures lighten up the thesis and simplify the explanations.
There are multiple ways to include a picture.
Which one fits best for your needs depends on the original image and the targeted quality of the print.
Overall, pixel based graphics are more frequent but results in bigger PDFs and a mediocre printing quality.
Vector based graphics are smaller and looks brilliant with a professional printer.
Nevertheless, it is much harder to create them.
The inclusion of both formats in LaTeX is quite similar.

\subsection{Drawing with \href{https://ctan.org/pkg/pgf}{\texttt{tikz}}}

This is the way you want to draw pictures in LaTeX!
But be warned, tikz's syntax is similarly powerful (or ugly) as plain LaTeX, and of course, you can use all LaTeX constructs in there as well.
This template is designed to have separate LaTeX files for each tikz picture, that can be compiled standalone.
That is very handy while drawing them because it allows real-time previews of the image using a suitable editor (e.g. gummi).
The final document includes the stored code like a regular graphic.
\Cref{fig:bluecircle} gives an example.

\begin{figure}[htb]
\centering
\documentclass[tikz]{standalone}
\begin{document}
\begin{tikzpicture}
  \fill [blue] (0,0) circle (1cm);
\end{tikzpicture}
\end{document}

% For a speed up with standalone tikz-pictures replace the above line with:
%\includestandalone{\myrootdir/img/circle}
\caption{A blue circle}
\label{fig:bluecircle}
\end{figure}

\subsection{Standard include}

OK, normal pictures can be included as well \ldots your images should, of course, be located in the img-folder \ldots not like this template does in \cref{fig:selogo} \ldots

\begin{figure}[htb]
\centering
\includegraphics[width=0.5\textwidth]{\myrootdir/zyx/TUMinfo/logos/2015_Logo_TUM_CMYK.pdf}
\caption{TUM logo scaled to half of the page width}
\label{fig:selogo}
\end{figure}

\newpage

\subsection{Side-by-side images}

If you run out of space, it might be helpful to put the images side by side as used as an example with \cref{fig:left-example} and \cref{fig:right-example}.

\begin{figure}[htb]
\centering
\begin{minipage}{0.46\textwidth}
\includegraphics[width=\textwidth]{example-image-a}
\caption{Left image}
\label{fig:left-example}
\end{minipage}
\qquad
\begin{minipage}{0.46\textwidth}
\includegraphics[width=\textwidth]{example-image-b}
\caption{Right image}
\label{fig:right-example}
\end{minipage}
\end{figure}


\section{Presenting data}
\label{sec:presenting-data}

LaTeX provides many ways to present your data in beautiful ways.
This template shows only a few of them, but hopefully, they are the most useful ones.
Most importantly, this \texttt{.tex} does not contain a single data point shown in the following plots and tables.
Everything is read from corresponding files in the \texttt{data}-directory.
I highly recommend keeping this approach for your visualizations as well, as these external files can easily be generated by external programs and can be updated flawlessly if any data changes.

\subsection{Beatiful plots}

\begin{figure}
\centering
\providecommand{\myrootdir}{..}
\documentclass[tikz]{standalone}
\onlyifstandalone{
\usepackage{pgfplots}
% This small hack is required to set the search path correctly
% Just provide the current directory this file is in relatively to the root dir
\pgfplotsset{table/search path={./img/}}
}
\begin{document}
\begin{tikzpicture}
 \begin{axis}[
   xlabel={input $i$},
   ylabel={result},
   grid=major,
   legend pos=north west,
  ]
  \addlegendentry{$i+2$}
  \addplot[blue,mark=x,thick] table[y index=1] {\myrootdir/data/plot-mult.dat};
  \addlegendentry{$2 \cdot i$}
  \addplot[red,mark=x,thick] table[y index=2] {\myrootdir/data/plot-mult.dat};
  \addlegendentry{$-1^{i}$}
  \addplot[green,mark=x,thick] table[y index=3] {\myrootdir/data/plot-mult.dat};
\end{axis}
\end{tikzpicture}
\end{document}

% For a speed up with standalone tikz-pictures replace the above line with:
%\includestandalone{\myrootdir/img/plot}
\caption{Example plot from external \texttt{.dat}-file}
\label{fig:plot-mult}
\end{figure}

When presenting a lot of complex data, visualizing them with plots is one of the best approaches you can choose.
But don't forget to separate data from code.
By the way, \texttt{pgfplot} can do way more things than you can see in \cref{fig:plot-mult}.

\subsection{Beautiful tables}

There are dozens of packages in LaTeX to compose tables.
This template focuses on the \href{https://ctan.org/pkg/booktabs}{\texttt{booktabs}}-package, but you can use what ever table fits best for you.
For using getting the best control over the inclusion of external \texttt{csv}-files you should familiarize yourself with the \href{https://www.ctan.org/pkg/csvsimple}{\texttt{csvsimple}}-package.

\subsubsection{Tables with a lot of numbers}
When you have tons of numbers, why not use a table filling a complete page like \cref{tab:101}

\begin{table}[p]
\begin{tabular*}{\textwidth}{@{\extracolsep{\fill$\;$}}c|*{10}{r}@{}}
\toprule
  \headcol{$\times$} & 1 & 2 & 3 & 4 & 5 & 6 & 7 & 8 & 9 & 10 \\
\midrule
\csvreader[no head,late after line=\\,late after last line=\\\bottomrule]{\myrootdir/data/101.csv}{}{%
\thecsvrow & \csvlinetotablerow}%
\end{tabular*}
\caption{The small $1 \cdot 1$}
\label{tab:101}
\end{table}

\clearpage % Small trick for better page breaks
\phantom{.}\vspace{-4em} % Big trick for better page breaks ^^
\subsubsection{Better number alignments}

\begin{wraptable}{o}{0.2\textwidth}
\vspace{-4em} % This number must be adjusted when used in a different context
\centering
\begin{tabular}{rD{.}{.}{3}}
\toprule
 \headcol{$n$} & \headcol{$n \cdot \sqrt{n}$} \\
\midrule
\csvreader[no head,late after line=\\,late after last line=\\\bottomrule]{\myrootdir/data/squares.csv}{}{%
\thecsvrow & \csvcolii}%
\end{tabular}
\caption{Square~roots}
\label{tab:roots}
\vspace{-2.5em} % This number must be adjusted when used in a different context
\end{wraptable}

When showing a lot of numbers with decimals, LaTeX offers a special alignment on the dot between the digits, what looks much better than just align right (see \cref{tab:roots}).
Have a look at the \href{https://ctan.org/pkg/dcolumn}{\texttt{dcolumn}}-package for more details.

This table also shows another neat feature:
You can put small tables on the side of your page and let the text float around it.
Of course, there is not much space for the caption, but for some tiny tables, it might also be helpful.

\subsubsection{Side-by-side tables}

If your tables are only very narrow, why not split them in two and show them more space efficiently side by side as in \cref{tab:twocolumn}?

\begin{table}[htb]
\begin{tabular*}{0.48\textwidth}{@{\extracolsep{\fill$\;$}}ccc@{}}
\toprule
  \headcol{Alpha} & \headcol{Beta} & \headcol{Gamma} \\
\midrule
\csvreader[head to column names,late after line=\\,late after last line=\\\bottomrule,%
% The trick is the filter expression with the line numer
filter expr={test{\ifnumless{\thecsvinputline}{5}}}%
]{\myrootdir/data/sidebyside.csv}{}{\Alpha & \Beta & \Gamma}
\end{tabular*}
%%%%%
\hfill
\begin{tabular*}{0.48\textwidth}{@{\extracolsep{\fill$\;$}}ccc@{}}
\toprule
  \headcol{Alpha} & \headcol{Beta} & \headcol{Gamma} \\
\midrule
\csvreader[head to column names,late after line=\\,late after last line=\\\bottomrule,%
% The trick is the filter expression with the line numer
filter expr={test{\ifnumgreater{\thecsvinputline}{4}}}%
]{\myrootdir/data/sidebyside.csv}{}{\Alpha & \Beta & \Gamma}
\end{tabular*}
\caption{Example for a table spread over two sides}
\label{tab:twocolumn}
\end{table}

\vspace{-1.5em} % Small trick for better page breaks
\section{Printing source code}

Including source code is simple and may look quite beautiful.
I recommend putting the actual code as a distinct file into the code directory with the correct extension.
Thereby you get perfect syntax highlighting from your local editor, and you can even compile or execute these files separately.
Latex can include each file with a single line of code.
You only must not use too many characters in one line in your code to avoid overlaps.

\lstinputlisting[float=t,language=C,caption=Hello World,label=lst:hello]{\myrootdir/code/hello.c}


\section{Miscellaneous}

While writing, you may leave some open tasks as todos for later.
Of course, there is also a command for this: \todo{Search for TODO in the text and fix them}
Furthermore, you can use all types of symbols (e.g. $\diameter$, €, \ldots).
And finally just a small citation from a random book by \textcite{Kemper2015} -- just to add the bibliography.
If you need direct quotes, just use the \enquote{quote-command} for simple quotation.

\end{document}
